\bookchapter{נספח א': תהליך “ברמת מטפלת” לפתיחת קופסת טלפון חדש}

\booksection{“טרילוגיית” בדיקת המכשיר עם הפתיחה:}

\begin{itemize}
  \item \textbf{בדיקת פרטים}: ודאו שה-IMEI/מספר סידורי על הקופסה תואם לחשבונית.
  \item \textbf{בדיקת פונקציות}: חייגו לבדיקת אפרכסת ומיקרופון, צלמו תמונה לבדיקת המצלמה.
  \item \textbf{2. "תזכורות חשובות" להפעלה והגדרה ראשונית:}רשת
\end{itemize}

\booksection{“תזכורות חשובות” להפעלה והגדרה ראשונית:}

\begin{itemize}
  \item \textbf{חשבון}: היכנסו ל-Apple ID, חשבון Huawei, חשבון Xiaomi וכו'.
  \item \textbf{הגדרות בסיסיות}: הפעילו את "מצא את הטלפון שלי", הגדירו סיסמת נעילה אמינה.
  \item \textbf{הגדרת eSIM}: אם הטלפון שלכם תומך ב-eSIM (כמו iPhone, חלק ממכשירי הדגל של אנדרואיד), תוכלו להוסיף חבילת eSIM בהגדרות, ללא צורך בכרטיס SIM פיזי. eSIM נוח במיוחד בנסיעות לחו"ל – ניתן לרכוש ישירות חבילת eSIM ממפעיל מקומי.
  \item \textbf{העברת נתונים}: בעת העברת נתונים מטלפון ישן, ודאו ששני הטלפונים טעונים מספיק. iPhone – השתמשו ב"התחלה מהירה", אנדרואיד – השתמשו באפליקציית המעבר הרשמית.
  \item \textbf{3. התקנת אפליקציות "בטיחות תחילה":}העדיפו חנות רשמית
\end{itemize}

\booksection{התקנת אפליקציות “בטיחות תחילה”:}

\begin{itemize}
  \item \textbf{אל תתנו הרשאות סתם}: כשקופצת בקשות הרשאות בהתקנה, חשבו "האם האפליקציה באמת צריכה את זה?"
  \item \textbf{נספח ב': רפרנס קצר לתרגילים מעשיים}תרגילי יסוד
\end{itemize}
